\documentclass[11pt]{article}
\usepackage[a4paper,margin=0.65in]{geometry}
\usepackage[T1]{fontenc}
\usepackage{amsmath,amssymb,amsfonts}
\usepackage{graphicx}
\usepackage{pdfpages}
\usepackage[english,shorthands=off]{babel}
\usepackage{fancyhdr}
\usepackage[bookmarks=false,colorlinks=false]{hyperref}
\usepackage[protrusion=true,expansion=false]{microtype}
\emergencystretch=3em
\tolerance=600
\providecommand{\modd}{\mathrm{mod}}
\providecommand{\Disc}{\mathrm{Disc}}
\providecommand{\canont}{}
\providecommand{\asterism}{*}
\providecommand{\Hessian}{H}
\providecommand{\tome}{}
\providecommand{\page}{p.}
\pagestyle{fancy}\fancyhf{}\fancyhead[L]{Pages 326--350}\fancyhead[R]{Cayley --- Collected Papers, Vol. XI}\renewcommand{\headrulewidth}{0.4pt}
\setlength{\headheight}{15pt}

\begin{document}

\begin{center}
{\small 302\hfill ON THE THEOREMS OF THE 2, 4, 8, AND 16 SQUARES.\hfill [763}
\end{center}

\begin{center}\textbf{II.}\end{center}

\begin{center}
\begin{tabular}{cccc|cccc|cccc|cccc}
1 & 2 & 3 & 4 & 5 & 6 & 7 & 8 & 9 & 10 & 11 & 12 & 13 & 14 & 15 & 16 \\
2 & 1 & 4 & 3 & 6 & 5 & 8 & 7 & 10 & 9 & 12 & 11 & 14 & 13 & 16 & 15 \\
3 & 4 & 1 & 2 & 7 & 8 & 5 & 6 & 11 & 12 & 9 & 10 & 15 & 16 & 13 & 14 \\
4 & 3 & 2 & 1 & 8 & 7 & 6 & 5 & 12 & 11 & 10 & 9 & 16 & 15 & 14 & 13 \\
\hline
5 & 6 & 7 & 8 & 1 & 2 & 3 & 4 & 14 & 13 & 16 & 15 & 10 & 9 & 12 & 11 \\
6 & 5 & 8 & 7 & 2 & 1 & 4 & 3 & 13 & 14 & 15 & 16 & 9 & 10 & 11 & 12 \\
7 & 8 & 5 & 6 & 3 & 4 & 1 & 2 & 16 & 15 & 14 & 13 & 12 & 11 & 10 & 9 \\
8 & 7 & 6 & 5 & 4 & 3 & 2 & 1 & 15 & 16 & 13 & 14 & 11 & 12 & 9 & 10 \\
\hline
9 & 10 & 11 & 12 & 14 & 13 & 16 & 15 & 1 & 2 & 3 & 4 & 6 & 5 & 8 & 7 \\
10 & 9 & 12 & 11 & 13 & 14 & 15 & 16 & 2 & 1 & 4 & 3 & 5 & 6 & 7 & 8 \\
11 & 12 & 9 & 10 & 16 & 15 & 14 & 13 & 3 & 4 & 1 & 2 & 8 & 7 & 6 & 5 \\
12 & 11 & 10 & 9 & 15 & 16 & 13 & 14 & 4 & 3 & 2 & 1 & 7 & 8 & 5 & 6 \\
\hline
13 & 14 & 15 & 16 & 10 & 9 & 12 & 11 & 6 & 5 & 8 & 7 & 1 & 2 & 3 & 4 \\
14 & 13 & 16 & 15 & 9 & 10 & 11 & 12 & 5 & 6 & 7 & 8 & 2 & 1 & 4 & 3 \\
15 & 16 & 13 & 14 & 12 & 11 & 10 & 9 & 8 & 7 & 6 & 5 & 3 & 4 & 1 & 2 \\
16 & 15 & 14 & 13 & 11 & 12 & 9 & 10 & 7 & 8 & 5 & 6 & 4 & 3 & 2 & 1 \\
\end{tabular}
\end{center}

\begin{center}\textbf{III.}\end{center}

\begin{center}
\begin{tabular}{cccc|cccc|cccc|cccc}
1 & 2 & 3 & 4 & 5 & 6 & 7 & 8 & 9 & 10 & 11 & 12 & 13 & 14 & 15 & 16 \\
2 & 1 & 4 & 3 & 6 & 5 & 8 & 7 & 10 & 9 & 12 & 11 & 14 & 13 & 16 & 15 \\
3 & 4 & 1 & 2 & 7 & 8 & 5 & 6 & 11 & 12 & 9 & 10 & 15 & 16 & 13 & 14 \\
4 & 3 & 2 & 1 & 8 & 7 & 6 & 5 & 12 & 11 & 10 & 9 & 16 & 15 & 14 & 13 \\
\hline
5 & 6 & 7 & 8 & 1 & 2 & 3 & 4 & 15 & 16 & 13 & 14 & 11 & 12 & 9 & 10 \\
6 & 5 & 8 & 7 & 2 & 1 & 4 & 3 & 16 & 15 & 14 & 13 & 12 & 11 & 10 & 9 \\
7 & 8 & 5 & 6 & 3 & 4 & 1 & 2 & 13 & 14 & 15 & 16 & 9 & 10 & 11 & 12 \\
8 & 7 & 6 & 5 & 4 & 3 & 2 & 1 & 14 & 13 & 16 & 15 & 10 & 9 & 12 & 11 \\
\hline
9 & 10 & 11 & 12 & 15 & 16 & 13 & 14 & 1 & 2 & 3 & 4 & 7 & 8 & 5 & 6 \\
10 & 9 & 12 & 11 & 16 & 15 & 14 & 13 & 2 & 1 & 4 & 3 & 8 & 7 & 6 & 5 \\
11 & 12 & 9 & 10 & 13 & 14 & 15 & 16 & 3 & 4 & 1 & 2 & 5 & 6 & 7 & 8 \\
12 & 11 & 10 & 9 & 14 & 13 & 16 & 15 & 4 & 3 & 2 & 1 & 6 & 5 & 8 & 7 \\
\hline
13 & 14 & 15 & 16 & 11 & 12 & 9 & 10 & 7 & 8 & 5 & 6 & 1 & 2 & 3 & 4 \\
14 & 13 & 16 & 15 & 12 & 11 & 10 & 9 & 8 & 7 & 6 & 5 & 2 & 1 & 4 & 3 \\
15 & 16 & 13 & 14 & 9 & 10 & 11 & 12 & 5 & 6 & 7 & 8 & 3 & 4 & 1 & 2 \\
16 & 15 & 14 & 13 & 10 & 9 & 12 & 11 & 6 & 5 & 8 & 7 & 4 & 3 & 2 & 1 \\
\end{tabular}
\end{center}

\begin{center}
{\small 763]\hfill ON THE THEOREMS OF THE 2, 4, 8, AND 16 SQUARES.\hfill 303}
\end{center}

\begin{center}\textbf{IV.}\end{center}

\begin{center}
\begin{tabular}{cccc|cccc|cccc|cccc}
1 & 2 & 3 & 4 & 5 & 6 & 7 & 8 & 9 & 10 & 11 & 12 & 13 & 14 & 15 & 16 \\
2 & 1 & 4 & 3 & 6 & 5 & 8 & 7 & 10 & 9 & 12 & 11 & 14 & 13 & 16 & 15 \\
3 & 4 & 1 & 2 & 7 & 8 & 5 & 6 & 11 & 12 & 9 & 10 & 15 & 16 & 13 & 14 \\
4 & 3 & 2 & 1 & 8 & 7 & 6 & 5 & 12 & 11 & 10 & 9 & 16 & 15 & 14 & 13 \\
\hline
5 & 6 & 7 & 8 & 1 & 2 & 3 & 4 & 16 & 15 & 14 & 13 & 12 & 11 & 10 & 9 \\
6 & 5 & 8 & 7 & 2 & 1 & 4 & 3 & 15 & 16 & 13 & 14 & 11 & 12 & 9 & 10 \\
7 & 8 & 5 & 6 & 3 & 4 & 1 & 2 & 14 & 13 & 16 & 15 & 10 & 9 & 12 & 11 \\
8 & 7 & 6 & 5 & 4 & 3 & 2 & 1 & 13 & 14 & 15 & 16 & 9 & 10 & 11 & 12 \\
\hline
9 & 10 & 11 & 12 & 16 & 15 & 14 & 13 & 1 & 2 & 3 & 4 & 8 & 7 & 6 & 5 \\
10 & 9 & 12 & 11 & 15 & 16 & 13 & 14 & 2 & 1 & 4 & 3 & 7 & 8 & 5 & 6 \\
11 & 12 & 9 & 10 & 14 & 13 & 16 & 15 & 3 & 4 & 1 & 2 & 6 & 5 & 8 & 7 \\
12 & 11 & 10 & 9 & 13 & 14 & 15 & 16 & 4 & 3 & 2 & 1 & 5 & 6 & 7 & 8 \\
\hline
13 & 14 & 15 & 16 & 12 & 11 & 10 & 9 & 8 & 7 & 6 & 5 & 1 & 2 & 3 & 4 \\
14 & 13 & 16 & 15 & 11 & 12 & 9 & 10 & 7 & 8 & 5 & 6 & 2 & 1 & 4 & 3 \\
15 & 16 & 13 & 14 & 10 & 9 & 12 & 11 & 6 & 5 & 8 & 7 & 3 & 4 & 1 & 2 \\
16 & 15 & 14 & 13 & 9 & 10 & 11 & 12 & 5 & 6 & 7 & 8 & 4 & 3 & 2 & 1 \\
\end{tabular}
\end{center}

The foregoing investigation of the synthematic arrangements is exhaustive: it
thereby appears that there are at most four types, viz.\ that every synthematic
arrangement is of the type of one or other of the four arrangements above written
down. The real nature of these is perhaps more clearly seen by means of the
corresponding squares; and it will be observed, that there is in the first square a
repetition of parts without transposition, which does not occur in the other three
squares; this seems to suggest, that while the first square (and therefore the first
synthematic arrangement) is really of a distinct type, the other three squares (or syn-thematic
arrangements) may possibly belong to one and the same type. If this were
so, it would be sufficient to prove the 16-theorem (viz.\ the non-existence of the
16-square formula) for the first and for any one of the other three synthematic
arrangements; but I provisionally assume that the four types are really distinct, and
propose therefore to prove the theorem for each of the four arrangements separately.

The process is the same as for the 8-theorem; we require the tetrads 1234, \&c.,
contained in the synthematic arrangements. In any one of these, each line gives
$\tfrac{1}{2}8.7$, $=28$ tetrads, and the 15 lines give therefore $15.28$, $=420$ tetrads:
but we thus obtain each tetrad 3 times, or the number of the tetrads is $420\div 3$, $=140$.

For the four arrangements respectively, these are as follows: the word ``same''
means same as in column I.

\begin{center}
{\small 304\hfill ON THE THEOREMS OF THE 2, 4, 8, AND 16 SQUARES.\hfill [763}
\end{center}

\begin{center}
\begin{tabular}{l l l l l}
& I. & II. & III. & IV. \\
\hline
1\quad 2\quad 3\quad 4 & & & & \\
& 5\quad 6 & & & \\
& 7\quad 8 & & & \\
& 9\quad 10 & & & \\
& 11\quad 12 & & & \\
& 13\quad 14 & & & \\
& 15\quad 16 & & & \\[1ex]
1\quad 3 & 5\quad 7 & & & \\
& 6\quad 8 & & & \\
& 9\quad 11 & & & \\
& 10\quad 12 & & & \\
& 13\quad 15 & & & \\
& 14\quad 16 & & & \\[1ex]
1\quad 4 & 5\quad 8 & & & \\
& 6\quad 7 & & & \\
& 9\quad 12 & & & \\
& 10\quad 11 & & & \\
& 13\quad 16 & & & \\
& 14\quad 15 & & & \\[1ex]
1\quad 5 & 9\quad 13 & 9\quad 14 & 9\quad 15 & 9\quad 16 \\
& 10\quad 14 & 10\quad 13 & 10\quad 16 & 10\quad 15 \\
& 11\quad 15 & 11\quad 16 & 11\quad 13 & 11\quad 14 \\
& 12\quad 16 & 12\quad 15 & 12\quad 14 & 12\quad 13 \\[1ex]
1\quad 6 & 9\quad 14 & 9\quad 13 & 9\quad 16 & 9\quad 15 \\
& 10\quad 13 & 10\quad 14 & 10\quad 15 & 10\quad 16 \\
& 11\quad 16 & 11\quad 15 & 11\quad 14 & 11\quad 13 \\
& 12\quad 15 & 12\quad 16 & 12\quad 13 & 12\quad 14 \\[1ex]
1\quad 7 & 9\quad 15 & 9\quad 16 & 9\quad 13 & 9\quad 14 \\
& 10\quad 16 & 10\quad 15 & 10\quad 14 & 10\quad 13 \\
& 11\quad 13 & 11\quad 14 & 11\quad 15 & 11\quad 16 \\
& 12\quad 14 & 12\quad 13 & 12\quad 16 & 12\quad 15 \\[1ex]
1\quad 8 & 9\quad 16 & 9\quad 15 & 9\quad 14 & 9\quad 13 \\
& 10\quad 15 & 10\quad 16 & 10\quad 13 & 10\quad 14 \\
& 11\quad 14 & 11\quad 13 & 11\quad 16 & 11\quad 15 \\
& 12\quad 13 & 12\quad 14 & 12\quad 15 & 12\quad 16 \\
\end{tabular}
\end{center}

\begin{center}
{\small 763]\hfill ON THE THEOREMS OF THE 2, 4, 8, AND 16 SQUARES.\hfill 305}
\end{center}

\begin{center}
\begin{tabular}{l l l l l}
& I. & II. & III. & IV. \\
\hline
2\quad 3 & 5\quad 8 & same & same & same \\
& 6\quad 7 & & & \\
& 9\quad 12 & & & \\
& 10\quad 11 & & & \\
& 13\quad 16 & & & \\
& 14\quad 15 & & & \\[1ex]
2\quad 4 & 5\quad 7 & same & same & same \\
& 6\quad 8 & & & \\
& 9\quad 11 & & & \\
& 10\quad 12 & & & \\
& 13\quad 15 & & & \\
& 14\quad 16 & & & \\[1ex]
2\quad 5 & 9\quad 14 & 9\quad 13 & 9\quad 16 & 9\quad 15 \\
& 10\quad 13 & 10\quad 14 & 10\quad 15 & 10\quad 16 \\
& 11\quad 16 & 11\quad 15 & 11\quad 14 & 11\quad 13 \\
& 12\quad 15 & 12\quad 16 & 12\quad 13 & 12\quad 14 \\[1ex]
2\quad 6 & 9\quad 13 & 9\quad 14 & 9\quad 15 & 9\quad 16 \\
& 10\quad 14 & 10\quad 13 & 10\quad 16 & 10\quad 15 \\
& 11\quad 15 & 11\quad 16 & 11\quad 13 & 11\quad 14 \\
& 12\quad 16 & 12\quad 15 & 12\quad 14 & 12\quad 13 \\[1ex]
2\quad 7 & 9\quad 16 & 9\quad 15 & 9\quad 14 & 9\quad 13 \\
& 10\quad 15 & 10\quad 16 & 10\quad 13 & 10\quad 14 \\
& 11\quad 14 & 11\quad 13 & 11\quad 16 & 11\quad 15 \\
& 12\quad 13 & 12\quad 14 & 12\quad 15 & 12\quad 16 \\[1ex]
2\quad 8 & 9\quad 15 & 9\quad 16 & 9\quad 13 & 9\quad 14 \\
& 10\quad 16 & 10\quad 15 & 10\quad 14 & 10\quad 13 \\
& 11\quad 13 & 11\quad 14 & 11\quad 15 & 11\quad 16 \\
& 12\quad 14 & 12\quad 13 & 12\quad 16 & 12\quad 15 \\[1ex]
3\quad 4 & 5\quad 6 & same & same & same \\
& 7\quad 8 & & & \\
& 9\quad 10 & & & \\
& 11\quad 12 & & & \\
& 13\quad 14 & & & \\
& 15\quad 16 & & & \\
\end{tabular}
\end{center}

\begin{center}
{\small C. XI.\hfill 39}
\end{center}

\begin{center}
{\small 306\hfill ON THE THEOREMS OF THE 2, 4, 8, AND 16 SQUARES.\hfill [763}
\end{center}

\begin{center}
\begin{tabular}{l l l l l}
& I. & II. & III. & IV. \\
\hline
3\quad 5 & 9\quad 15 & 9\quad 16 & 9\quad 13 & 9\quad 14 \\
& 10\quad 16 & 10\quad 15 & 10\quad 14 & 10\quad 13 \\
& 11\quad 13 & 11\quad 14 & 11\quad 15 & 11\quad 16 \\
& 12\quad 14 & 12\quad 13 & 12\quad 16 & 12\quad 15 \\[1ex]
3\quad 6 & 9\quad 16 & 9\quad 15 & 9\quad 14 & 9\quad 13 \\
& 10\quad 15 & 10\quad 16 & 10\quad 13 & 10\quad 14 \\
& 11\quad 14 & 11\quad 13 & 11\quad 16 & 11\quad 15 \\
& 12\quad 13 & 12\quad 14 & 12\quad 15 & 12\quad 16 \\[1ex]
3\quad 7 & 9\quad 13 & 9\quad 14 & 9\quad 15 & 9\quad 16 \\
& 10\quad 14 & 10\quad 13 & 10\quad 16 & 10\quad 15 \\
& 11\quad 15 & 11\quad 16 & 11\quad 13 & 11\quad 14 \\
& 12\quad 16 & 12\quad 15 & 12\quad 14 & 12\quad 13 \\[1ex]
3\quad 8 & 9\quad 14 & 9\quad 13 & 9\quad 16 & 9\quad 15 \\
& 10\quad 13 & 10\quad 14 & 10\quad 15 & 10\quad 16 \\
& 11\quad 16 & 11\quad 15 & 11\quad 14 & 11\quad 13 \\
& 12\quad 15 & 12\quad 16 & 12\quad 13 & 12\quad 14 \\[1ex]
4\quad 5 & 9\quad 16 & 9\quad 15 & 9\quad 14 & 9\quad 13 \\
& 10\quad 15 & 10\quad 16 & 10\quad 13 & 10\quad 14 \\
& 11\quad 14 & 11\quad 13 & 11\quad 16 & 11\quad 15 \\
& 12\quad 13 & 12\quad 14 & 12\quad 15 & 12\quad 16 \\[1ex]
4\quad 6 & 9\quad 15 & 9\quad 16 & 9\quad 13 & 9\quad 14 \\
& 10\quad 16 & 10\quad 15 & 10\quad 14 & 10\quad 13 \\
& 11\quad 13 & 11\quad 14 & 11\quad 15 & 11\quad 16 \\
& 12\quad 14 & 12\quad 13 & 12\quad 16 & 12\quad 15 \\[1ex]
4\quad 7 & 9\quad 14 & 9\quad 13 & 9\quad 16 & 9\quad 15 \\
& 10\quad 13 & 10\quad 14 & 10\quad 15 & 10\quad 16 \\
& 11\quad 16 & 11\quad 15 & 11\quad 14 & 11\quad 13 \\
& 12\quad 15 & 12\quad 16 & 12\quad 13 & 12\quad 14 \\[1ex]
4\quad 8 & 9\quad 13 & 9\quad 14 & 9\quad 15 & 9\quad 16 \\
& 10\quad 14 & 10\quad 13 & 10\quad 16 & 10\quad 15 \\
& 11\quad 15 & 11\quad 16 & 11\quad 13 & 11\quad 14 \\
& 12\quad 16 & 12\quad 15 & 12\quad 14 & 12\quad 13 \\[1ex]
5\quad 6 & 7\quad 8 & same & same & same \\
& 9\quad 10 & & & \\
& 11\quad 12 & & & \\
& 13\quad 14 & & & \\
& 15\quad 16 & & & \\
\end{tabular}
\end{center}

\begin{center}
{\small 763]\hfill ON THE THEOREMS OF THE 2, 4, 8, AND 16 SQUARES.\hfill 307}
\end{center}

\begin{center}
\begin{tabular}{l l l l l}
& I. & II. & III. & IV. \\
\hline
5\quad 7 & 9\quad 11 & same & same & same \\
& 10\quad 12 & & & \\
& 13\quad 15 & & & \\
& 14\quad 16 & & & \\[1ex]
5\quad 8 & 9\quad 12 & & & \\
& 10\quad 11 & & & \\
& 13\quad 16 & & & \\
& 14\quad 15 & & & \\[1ex]
6\quad 7 & 9\quad 12 & & & \\
& 10\quad 11 & & & \\
& 13\quad 16 & & & \\
& 14\quad 15 & & & \\[1ex]
6\quad 8 & 9\quad 11 & & & \\
& 10\quad 12 & & & \\
& 13\quad 15 & & & \\
& 14\quad 16 & & & \\[1ex]
7\quad 8 & 9\quad 10 & & & \\
& 11\quad 12 & & & \\
& 13\quad 14 & & & \\
& 15\quad 16 & & & \\[1ex]
9\quad 10 & 11\quad 12 & & & \\
& 13\quad 14 & & & \\
& 15\quad 16 & & & \\[1ex]
9\quad 11 & 13\quad 15 & & & \\
& 14\quad 16 & & & \\[1ex]
9\quad 12 & 13\quad 16 & & & \\
& 14\quad 15 & & & \\[1ex]
10\quad 11 & 13\quad 16 & & & \\
& 14\quad 15 & & & \\[1ex]
10\quad 12 & 13\quad 15 & & & \\
& 14\quad 16 & & & \\[1ex]
11\quad 12 & 13\quad 14 & & & \\
& 15\quad 16 & & & \\[1ex]
13\quad 14 & 15\quad 16 & & & \\
\end{tabular}
\end{center}

\begin{center}{\small 39---2}\end{center}

\begin{center}
{\small 308\hfill ON THE THEOREMS OF THE 2, 4, 8, AND 16 SQUARES.\hfill [763}
\end{center}

As regards the signs, observe that the first line may always be written
\[
ab + cd + ef + \&c.,
\]
with the signs all of them $+$; and then writing $a, b, c, \ldots = 1, 2, 3, \ldots, 16$ respectively,
the first line will be
\[
1\;2 + 3\;4 + 5\;6 + 7\;8 + 9\;10 + 11\;12 + 13\;14 + 15\;16,
\]
with the signs all of them $+$; that is, we may assume $e_{12}, e_{34},$ \&c., or say
\[
1\;2,\ 3\;4,\ 5\;6,\ 7\;8,\ 9\;10,\ 11\;12,\ 13\;14,\ 15\;16,
\]
all of them $=+1$. And in the other lines, the signs of all the terms of any line
may be reversed at pleasure, that is, we may assume $e_{13}, e_{14},$ \&c., or say $1\;3, 1\;4,
1\;5, 1\;6, 1\;7, 1\;8, 1\;9, 1\;10, 1\;11, 1\;12, 1\;13, 1\;14, 1\;15, 1\;16$, all of them
$=+1$.

Making these assumptions, then for any one of the synthematic arrangements the
several tetrads give as before relations between the signs; among these are included
the results already obtained for the 8-question, and taking as before
\[
-a = b = c = d = \theta,
\]
we have the signs of the several terms belonging to the 8-question given as $= +\theta$
or $+\beta$, as before. The remaining tetrads up to $1\;8\;12\;13$ then serve to express all
the remaining signs in terms of the as yet undetermined signs $e, f, g, h, i, j, k, l,
m, n, o, p, q, r, s, t, u, v, w, x, y, z, \alpha, \beta$, for instance
\[
1\;3.\,9\;11 + 1.\,e,
\]
\[
1\;9.\,3\;11 - 1.\,e,
\]
\[
1\;11.\,3\;9 + 1.\,e,
\]
that is, $3\;9 = e,\ 3\;11 = -e,\ 9\;11 = e$; and then the tetrads up to $2\;8\;9\;15$ serve to
express these signs in terms of the undetermined signs $\lambda, \mu, \nu, \rho, \sigma, \tau$; for instance
\[
2\;3.\,9\;12 + 1.\,i,
\]
\[
2\;9.\,3\;12 - 1.\,f,
\]
\[
2\;12.\,3\;9 + -1.\,e,
\]
that is, $-e = f = i$; and in like manner $2\;3\;10\;11$, $2\;4\;9\;11$ and $2\;4\;10\;12$ give
respectively $-e = f = j$, $-e = i = j$, $f = i = j$; that is, we have $-e = f = i = j$, $= \lambda$ suppose.
And in this way we have, for each of the four synthematic arrangements the signs
of all the terms expressed in terms of the undetermined signs $\theta, \lambda, \mu, \nu, \rho, \sigma, \tau$,
as shown in the following table; where observe that the results apply to the four
synthematic arrangements separately, viz.\ the $e, f, g, \ldots$ \&c., and the $\theta, \lambda, \mu, \nu, \rho, \sigma, \tau$
in each column are altogether independent of the like symbols in the other three
columns.

Signs for the four synthematic arrangements:

\begin{center}
{\small 763]\hfill ON THE THEOREMS OF THE 2, 4, 8, AND 16 SQUARES.\hfill 309}
\end{center}

\begin{center}
\begin{tabular}{l l l l l}
& I. & II. & III. & IV. \\
\hline
1\quad 2 & $+1$ & same & same & same \\
\phantom{1}\quad 3 & $+1$ & & & \\
\phantom{1}\quad 4 & $+1$ & & & \\
\phantom{1}\quad 5 & $+1$ & & & \\
\phantom{1}\quad 6 & $+1$ & & & \\
\phantom{1}\quad 7 & $+1$ & & & \\
\phantom{1}\quad 8 & $+1$ & & & \\
\phantom{1}\quad 9 & $+1$ & & & \\
\phantom{1}\quad 10 & $+1$ & & & \\
\phantom{1}\quad 11 & $+1$ & & & \\
\phantom{1}\quad 12 & $+1$ & & & \\
\phantom{1}\quad 13 & $+1$ & & & \\
\phantom{1}\quad 14 & $+1$ & & & \\
\phantom{1}\quad 15 & $+1$ & & & \\
\phantom{1}\quad 16 & $+1$ & & & \\[1ex]
2\quad 3 & $+1$ & & & \\
\phantom{2}\quad 4 & $-1$ & & & \\
\phantom{2}\quad 5 & $+1$ & & & \\
\phantom{2}\quad 6 & $-1$ & & & \\
\phantom{2}\quad 7 & $+1$ & & & \\
\phantom{2}\quad 8 & $-1$ & & & \\
\phantom{2}\quad 9 & $+1$ & & & \\
\phantom{2}\quad 10 & $-1$ & & & \\
\phantom{2}\quad 11 & $+1$ & & & \\
\phantom{2}\quad 12 & $-1$ & & & \\
\phantom{2}\quad 13 & $+1$ & & & \\
\phantom{2}\quad 14 & $-1$ & & & \\
\phantom{2}\quad 15 & $+1$ & & & \\
\phantom{2}\quad 16 & $-1$ & & & \\[1ex]
3\quad 4 & $+1$ & & same & same \\
\phantom{3}\quad 5 & $-\theta$ & & & \\
\phantom{3}\quad 6 & $\theta$ & & & \\
\phantom{3}\quad 7 & $\theta$ & & & \\
\phantom{3}\quad 8 & $-\theta$ & & & \\
\phantom{3}\quad 9 & $e\ -\lambda$ & & & \\
\phantom{3}\quad 10 & $f\ \lambda$ & & & \\
\phantom{3}\quad 11 & $-e\ \lambda$ & & & \\
\phantom{3}\quad 12 & $-f\ -\lambda$ & & & \\
\phantom{3}\quad 13 & $g\ -\mu$ & & & \\
\phantom{3}\quad 14 & $h\ \mu$ & & & \\
\phantom{3}\quad 15 & $-g\ \mu$ & & & \\
\phantom{3}\quad 16 & $-h\ -\mu$ & & & \\
\end{tabular}
\end{center}

\begin{center}
{\small 310\hfill ON THE THEOREMS OF THE 2, 4, 8, AND 16 SQUARES.\hfill [763}
\end{center}

\begin{center}
\begin{tabular}{l l l l l}
& I. & II. & III. & IV. \\
\hline
4\quad 5 & & same & same & same \\
\phantom{4}\quad 6 & $\theta$ & & & \\
\phantom{4}\quad 7 & $-\theta$ & & & \\
\phantom{4}\quad 8 & $-\theta$ & & & \\
\phantom{4}\quad 9 & $i\ \lambda$ & & & \\
\phantom{4}\quad 10 & $j\ \lambda$ & & & \\
\phantom{4}\quad 11 & $-j\ -\lambda$ & & & \\
\phantom{4}\quad 12 & $-i\ -\lambda$ & & & \\
\phantom{4}\quad 13 & $k\ \mu$ & & & \\
\phantom{4}\quad 14 & $l\ \mu$ & & & \\
\phantom{4}\quad 15 & $-l\ -\mu$ & & & \\
\phantom{4}\quad 16 & $-k\ -\mu$ & & & \\[1ex]
5\quad 6 & $+1$ & $+1$ & $+1$ & $+1$ \\
\phantom{5}\quad 7 & $-\theta$ & $-\theta$ & $-\theta$ & $-\theta$ \\
\phantom{5}\quad 8 & $\theta$ & & & \\
\phantom{5}\quad 9 & $m\ -\nu$ & $m\ \nu$ & $m\ -\nu$ & $m\ \nu$ \\
\phantom{5}\quad 10 & $n\ \nu$ & $n\ \nu$ & $n\ \nu$ & $n\ \nu$ \\
\phantom{5}\quad 11 & $o\ -\rho$ & $o\ \rho$ & $o\ -\rho$ & $o\ \rho$ \\
\phantom{5}\quad 12 & $p\ \rho$ & $p\ \rho$ & $p\ \rho$ & $p\ \rho$ \\
\phantom{5}\quad 13 & $-m\ \nu$ & $-n\ -\nu$ & $-o\ \rho$ & $-p\ -\rho$ \\
\phantom{5}\quad 14 & $-n\ -\nu$ & $-m\ -\nu$ & $-p\ -\rho$ & $-o\ -\rho$ \\
\phantom{5}\quad 15 & $-o\ \rho$ & $-p\ -\rho$ & $-m\ \nu$ & $-n\ -\nu$ \\
\phantom{5}\quad 16 & $-p\ -\rho$ & $-o\ -\rho$ & $-n\ -\nu$ & $-m\ -\nu$ \\[1ex]
6\quad 7 & & & $\theta$ & \\
\phantom{6}\quad 8 & $\theta$ & & & \\
\phantom{6}\quad 9 & $q\ \nu$ & $q\ \nu$ & $q\ \nu$ & $q\ \nu$ \\
\phantom{6}\quad 10 & $r\ \nu$ & $r\ -\nu$ & $r\ \nu$ & $r\ -\nu$ \\
\phantom{6}\quad 11 & $s\ \rho$ & $s\ \rho$ & $s\ \rho$ & $s\ \rho$ \\
\phantom{6}\quad 12 & $t\ \rho$ & $t\ -\rho$ & $t\ \rho$ & $t\ -\rho$ \\
\phantom{6}\quad 13 & $-r\ -\nu$ & $-q\ -\nu$ & $-t\ -\rho$ & $-s\ -\rho$ \\
\phantom{6}\quad 14 & $-q\ -\nu$ & $-r\ \nu$ & $-s\ -\rho$ & $-t\ \rho$ \\
\phantom{6}\quad 15 & $-t\ -\rho$ & $-s\ -\rho$ & $-r\ -\nu$ & $-q\ -\nu$ \\
\phantom{6}\quad 16 & $-s\ -\rho$ & $-t\ \rho$ & $-q\ -\nu$ & $-r\ \nu$ \\[1ex]
7\quad 8 & $+1$ & $+1$ & $+1$ & $+1$ \\
\phantom{7}\quad 9 & $u\ -\sigma$ & $u\ -\sigma$ & $u\ -\sigma$ & $u\ \sigma$ \\
\phantom{7}\quad 10 & $v\ \sigma$ & $v\ \sigma$ & $v\ \sigma$ & $v\ \sigma$ \\
\phantom{7}\quad 11 & $w\ -\tau$ & $w\ \tau$ & $w\ -\tau$ & $w\ \tau$ \\
\phantom{7}\quad 12 & $x\ \tau$ & $x\ \tau$ & $x\ \tau$ & $x\ \tau$ \\
\phantom{7}\quad 13 & $-w\ \tau$ & $-x\ -\tau$ & $-u\ \sigma$ & $-v\ -\sigma$ \\
\phantom{7}\quad 14 & $-x\ -\tau$ & $-w\ -\tau$ & $-v\ -\sigma$ & $-u\ -\sigma$ \\
\phantom{7}\quad 15 & $-u\ \sigma$ & $-v\ -\sigma$ & $-w\ \tau$ & $-x\ -\tau$ \\
\phantom{7}\quad 16 & $-v\ -\sigma$ & $-u\ -\sigma$ & $-x\ -\tau$ & $-w\ -\tau$ \\
\end{tabular}
\end{center}

\begin{center}
{\small 763]\hfill ON THE THEOREMS OF THE 2, 4, 8, AND 16 SQUARES.\hfill 311}
\end{center}

\begin{center}
\begin{tabular}{l l l l l}
& I. & II. & III. & IV. \\
\hline
8\quad 9 & $y\ \sigma$ & $y\ \sigma$ & $y\ \sigma$ & $y\ \sigma$ \\
\phantom{8}\quad 10 & $z\ \sigma$ & $z\ -\sigma$ & $z\ \sigma$ & $z\ -\sigma$ \\
\phantom{8}\quad 11 & $\alpha\ \tau$ & $\alpha\ \tau$ & $\alpha\ \tau$ & $\alpha\ \tau$ \\
\phantom{8}\quad 12 & $\beta\ \tau$ & $\beta\ -\tau$ & $\beta\ \tau$ & $\beta\ -\tau$ \\
\phantom{8}\quad 13 & $-\beta\ -\tau$ & $-\alpha\ -\tau$ & $-z\ -\sigma$ & $-y\ -\sigma$ \\
\phantom{8}\quad 14 & $-\alpha\ -\tau$ & $-\beta\ \tau$ & $-y\ -\sigma$ & $-z\ \sigma$ \\
\phantom{8}\quad 15 & $-z\ -\sigma$ & $-y\ -\sigma$ & $-\beta\ -\tau$ & $-\alpha\ -\tau$ \\
\phantom{8}\quad 16 & $-y\ -\sigma$ & $-z\ \sigma$ & $-\alpha\ -\tau$ & $-\beta\ \tau$ \\[1ex]
9\quad 10 & $+1$ & $+1$ & $+1$ & $+1$ \\
\phantom{9}\quad 11 & $e\ -\lambda$ & $e\ -\lambda$ & $e\ -\lambda$ & $e\ -\lambda$ \\
\phantom{9}\quad 12 & $i\ \lambda$ & $i\ \lambda$ & $i\ \lambda$ & $i\ \lambda$ \\
\phantom{9}\quad 13 & $m\ -\nu$ & $q\ \nu$ & $u\ -\sigma$ & $y\ \sigma$ \\
\phantom{9}\quad 14 & $q\ \nu$ & $m\ \nu$ & $y\ \sigma$ & $u\ \sigma$ \\
\phantom{9}\quad 15 & $u\ -\sigma$ & $y\ \sigma$ & $m\ -\nu$ & $q\ \nu$ \\
\phantom{9}\quad 16 & $y\ \sigma$ & $u\ \sigma$ & $q\ \nu$ & $m\ \nu$ \\[1ex]
10\quad 11 & $j\ \lambda$ & $j\ \lambda$ & $j\ \lambda$ & $j\ \lambda$ \\
\phantom{10}\quad 12 & $f\ \lambda$ & $f\ \lambda$ & $f\ \lambda$ & $f\ \lambda$ \\
\phantom{10}\quad 13 & $r\ \nu$ & $n\ \nu$ & $z\ \sigma$ & $v\ \sigma$ \\
\phantom{10}\quad 14 & $n\ \nu$ & $r\ -\nu$ & $v\ \sigma$ & $z\ -\sigma$ \\
\phantom{10}\quad 15 & $z\ \sigma$ & $v\ \sigma$ & $r\ \nu$ & $n\ \nu$ \\
\phantom{10}\quad 16 & $v\ \sigma$ & $z\ -\sigma$ & $n\ \nu$ & $r\ -\nu$ \\[1ex]
11\quad 12 & $+1$ & $+1$ & $+1$ & $+1$ \\
\phantom{11}\quad 13 & $w\ -\tau$ & $s\ \tau$ & $o\ -\rho$ & $s\ \rho$ \\
\phantom{11}\quad 14 & $s\ \tau$ & $w\ \tau$ & $s\ \rho$ & $o\ \rho$ \\
\phantom{11}\quad 15 & $o\ -\rho$ & $s\ \rho$ & $w\ -\tau$ & $s\ \tau$ \\
\phantom{11}\quad 16 & $s\ \rho$ & $o\ \rho$ & $s\ \tau$ & $w\ \tau$ \\[1ex]
12\quad 13 & $x\ \tau$ & $x\ \tau$ & $t\ \rho$ & $p\ \rho$ \\
\phantom{12}\quad 14 & $x\ \tau$ & $\beta\ -\tau$ & $p\ \rho$ & $t\ -\rho$ \\
\phantom{12}\quad 15 & $t\ \rho$ & $p\ \rho$ & $p\ \tau$ & $x\ \tau$ \\
\phantom{12}\quad 16 & $p\ \rho$ & $t\ -\rho$ & $x\ \tau$ & $p\ -\tau$ \\[1ex]
13\quad 14 & $+1$ & $+1$ & $+1$ & $+1$ \\
\phantom{13}\quad 15 & $g\ -\mu$ & $g\ -\mu$ & $g\ -\mu$ & $g\ -\mu$ \\
\phantom{13}\quad 16 & $k\ \mu$ & $k\ \mu$ & $k\ \mu$ & $k\ \mu$ \\[1ex]
14\quad 15 & $l\ \mu$ & $l\ \mu$ & $l\ \mu$ & $l\ \mu$ \\
\phantom{14}\quad 16 & $h\ \mu$ & $h\ \mu$ & $h\ \mu$ & $h\ \mu$ \\[1ex]
15\quad 16 & $+1$ & $+1$ & $+1$ & $+1$ \\
\end{tabular}
\end{center}

\begin{center}
{\small 312\hfill ON THE THEOREMS OF THE 2, 4, 8, AND 16 SQUARES.\hfill [763}
\end{center}

We have now for the four arrangements respectively, by means of hitherto unused
tetrads, the following determinations of sign: these being in each case inconsistent with
each other.

\textit{First arrangement.}

\begin{align*}
3\;5\;9\;15 & + -\theta.-\sigma & & \text{that is,} \\
3\;9\;5\;15 & - -\lambda.\ \rho & & \theta\sigma = -\lambda\rho = -\mu\nu, \\
3\;15\;5\;9 & + -\mu.-\nu
\end{align*}

\begin{align*}
3\;5\;10\;16 & + -\theta.\ \sigma & & \\
3\;10\;5\;16 & - -\lambda.-\rho & & -\theta\sigma = \lambda\rho = -\mu\nu. \\
3\;16\;5\;10 & + -\mu.\ \nu
\end{align*}

\begin{align*}
3\;5\;11\;13 & + -\theta.-\tau & & \\
3\;11\;5\;13 & - -\lambda.\ \nu & & \theta\tau = -\lambda\nu = \mu\rho, \\
3\;13\;5\;11 & + -\mu.-\rho
\end{align*}

\begin{align*}
3\;5\;12\;14 & + -\theta.\ \tau & & \\
3\;12\;5\;14 & - -\lambda.-\nu & & -\theta\tau = -\lambda\nu = \mu\rho. \\
3\;14\;5\;12 & + \ \mu.\ \rho
\end{align*}

\textit{Second arrangement.}

\begin{align*}
3\;5\;9\;16 & + -\theta.\ \sigma & & \\
3\;9\;5\;16 & - -\lambda.-\rho & & \theta\sigma = \lambda\rho = \mu\nu, \\
3\;16\;5\;9 & + -\mu.\ \nu
\end{align*}

\begin{align*}
3\;5\;10\;15 & + -\theta.\ \sigma & & \\
3\;10\;5\;15 & - -\lambda.-\rho & & -\theta\sigma = \lambda\rho = \mu\nu, \\
3\;15\;5\;10 & + \ \mu.\ \nu
\end{align*}

\begin{align*}
3\;5\;11\;14 & + -\theta.\ \tau & & \\
3\;11\;5\;14 & - -\lambda.-\nu & & -\theta\tau = \lambda\nu = \mu\rho, \\
3\;14\;5\;11 & + \ \mu.\ \rho
\end{align*}

\begin{align*}
3\;5\;12\;13 & + -\theta.\ \tau & & \\
3\;12\;5\;13 & - -\lambda.-\nu & & \theta\tau = \lambda\nu = \mu\rho. \\
3\;13\;5\;12 & + -\mu.\ \rho
\end{align*}

\begin{center}
{\small 763]\hfill ON THE THEOREMS OF THE 2, 4, 8, AND 16 SQUARES.\hfill 313}
\end{center}

\textit{Third arrangement.}

\begin{align*}
3\;5\;9\;13 & + -\theta.-\sigma & & \\
3\;9\;5\;13 & - -\lambda.\ \rho & & \theta\sigma = -\lambda\rho = \mu\nu, \\
3\;13\;5\;9 & + -\mu.-\nu
\end{align*}

\begin{align*}
3\;5\;10\;14 & + -\theta.\ \sigma & & \\
3\;10\;5\;14 & - -\lambda.-\rho & & \theta\sigma = -\lambda\rho = -\mu\nu, \\
3\;14\;5\;10 & + \ \mu.\ \nu
\end{align*}

\begin{align*}
3\;5\;11\;15 & + -\theta.-\tau & & \\
3\;11\;5\;15 & - -\lambda.\ \nu & & \theta\tau = -\lambda\nu = -\mu\rho. \\
3\;15\;5\;11 & + \ \mu.-\rho
\end{align*}

\begin{align*}
3\;5\;12\;16 & + -\theta.\ \tau & & \\
3\;12\;5\;16 & - -\lambda.-\nu & & \theta\tau = -\lambda\nu = \mu\rho. \\
3\;16\;5\;12 & + -\mu.\ \rho
\end{align*}

\textit{Fourth arrangement.}

\begin{align*}
3\;5\;9\;14 & + -\theta.\ \sigma & & \\
3\;9\;5\;14 & - -\lambda.-\rho & & \theta\sigma = \lambda\rho = -\mu\nu, \\
3\;14\;5\;9 & + \ \mu.\ \nu
\end{align*}

\begin{align*}
3\;5\;10\;13 & + -\theta.\ \sigma & & \\
3\;10\;5\;13 & - -\lambda.-\rho & & \theta\sigma = -\lambda\rho = \mu\nu, \\
3\;13\;5\;10 & + -\mu.\ \nu
\end{align*}

\begin{align*}
3\;5\;11\;16 & + -\theta.\ \tau & & \\
3\;11\;5\;16 & - -\lambda.-\nu & & \theta\tau = -\lambda\nu = \mu\rho, \\
3\;16\;5\;11 & + -\mu.\ \rho
\end{align*}

\begin{align*}
3\;5\;12\;15 & + -\theta.-\tau & & \\
3\;12\;5\;15 & - -\lambda.-\nu & & \theta\tau = -\lambda\nu = -\mu\rho. \\
3\;15\;5\;12 & + -\mu.\ \rho
\end{align*}

And it hence finally appears, that we cannot, in any one of the four arrange-ments,
determine the signs so as to give rise to a 16-square theorem; that is, the
product of a sum of 16 squares into a sum of 16 squares cannot be made equal to
a sum of 16 squares.

\begin{center}{\small C. XI.\hfill 40}\end{center}

\bigskip

\begin{center}
{\small 314\hfill\hfill [764}
\end{center}

\begin{center}\textbf{764.}\end{center}

\begin{center}\textbf{THE BINOMIAL EQUATION $x^n - 1 = 0$: QUINQUISECTION.}\end{center}

\begin{center}\small [From the \textit{Proceedings of the London Mathematical Society}, vol.\ \textsc{xii}.\ (1881), pp.\ 15, 16.\\
Read December 9, 1880.]\end{center}

\textsc{The} theory should be precisely analogous to those for the trisection and quarti-section
(see my paper, ``The Binomial Equation $x^n - 1 = 0$, Trisection and Quartisection,''
\textit{Proceedings of the London Mathematical Society}, vol.\ \textsc{xi}.\ (1879), pp.\ 4--17, [731]),
only I have not been able to carry it so far. We have in the present case five
periods $X, Y, Z, W, T$, the actual expressions for which, $X = \eta^a + \ldots$, $Y = \eta^b + \ldots$, etc.,
with Reuschle's selected prime root $g$, can be (for the primes $5n + 1$ under 100) at
once written down by means of the table given, pp.\ 16, 17, of that paper; [see this
volume, pp.\ 95, 96]. The relations between the periods are of the form
\[
\begin{array}{c|ccccc}
 & X & Y & Z & W & T \\
\hline
X^2 = & a & b & c & d & e \\
XY = & f & g & h & i & j \\
XZ = & k & l & m & n & o
\end{array}
\]
that is, we have
\[
X^2 = (a, b, c, d, e \mid X, Y, Z, W, T),
\]
and thence, by cyclical permutations,
\[
Y^2 = (e, a, b, c, d \mid \quad ),\ \text{etc.};
\]
viz.\ from the value of $X^2$ we have those of $Y^2, Z^2, W^2, T^2$; from the value of $XY$
those of $YZ, ZW, WT, TX$; and from the value of $XZ$ those of $YW, ZT, WX, TY$.

\begin{center}
{\small 764]\hfill THE BINOMIAL EQUATION $x^n - 1 = 0$: QUINQUISECTION.\hfill 315}
\end{center}

From the equation $X + Y + Z + W + T = -1$, multiplying by $X$ and then substi-tuting
for $X^2$, $XY$, \&c., their values, we obtain
\begin{align*}
-a &= 1 + f + k + m + g, \\
-b &= \phantom{1 +\,} g + l + n + h, \\
-c &= \phantom{1 +\,} h + m + o + i, \\
-d &= \phantom{1 +\,} i + n + k + j, \\
-e &= \phantom{1 +\,} j + o + l + f,
\end{align*}
which determine $(a, b, c, d, e)$ in terms of $(f, g, h, i, j)$ and $(k, l, m, n, o)$. It is,
moreover, easy to prove that
\begin{align*}
f + g + h + i + j &= \tfrac{1}{5}(p - 1), \\
k + l + m + n + o &= \tfrac{1}{5}(p - 1),
\end{align*}
whence also
\[
a + b + c + d + e = -1 - \tfrac{2}{5}(p - 1).
\]
We obtain other relations between the coefficients by considering the two triple
products $XYZ$ and $XYW$: these are all that need be considered, since the other
triple products are deducible from them by cyclical permutations. From the first of
these we have
\[
X.YZ = Y.XZ = Z.XY,
\]
and from the second
\[
X.YW = Y.XW = W.XY;
\]
and if we herein substitute for $YZ, XZ$, \&c., their values, and then in the resulting
equations for $X^2, XY$, \&c., their values as linear functions of $X, Y, Z, W, T$, we
obtain in all $5.2.2 = 20$ quadric relations between the 15 coefficients; or if we
substitute for $(a, b, c, d, e)$ their foregoing values, in all 20 relations between the 10
coefficients $(f, g, h, i, j)$ and $(k, l, m, n, o)$. These are at most equivalent to 8
independent equations, since we have, besides, the sums $f+g+h+i+j$ and $k+l+m+n+o$
each $= \tfrac{1}{5}(p - 1)$; but I have not succeeded in finding the connexions between them,
or even in ascertaining to how many independent equations they are equivalent.

For any given prime $p = 5n + 1$, the values of the coefficients, and also the
coefficients of the quintic equation for the periods, could of course be calculated
directly from the expressions of the periods; but for the primes under 100, that is,
for the values 11, 31, 41, 61, 71, they are at once obtained from Reuschle. We
have thus the two Tables, the former giving the coefficients $a, b, \ldots, n, o$, and the
latter the coefficients of the quintic equations.

\begin{center}{\small 40---2}\end{center}

\begin{center}
{\small 316\hfill THE BINOMIAL EQUATION $x^n - 1 = 0$: QUINQUISECTION.\hfill [764}
\end{center}

\begin{center}
\textsc{Table 1.}\hspace{2cm}\textsc{Table 2} of the Quintic Equations.
\end{center}

\begin{center}
\begin{tabular}{c|rrrrr}
& $a$ & $b$ & $c$ & $d$ & $e$ \\
$p$ & $f$ & $g$ & $h$ & $i$ & $j$ \\
& $k$ & $l$ & $m$ & $n$ & $o$ \\
\hline
11 & $-2$ & $-1$ & $-2$ & $-2$ & $-2$ \\
& 1 & 0 & 0 & 1 & 0 \\
& 0 & 0 & 0 & 1 & 1 \\
\hline
31 & $-4$ & $-6$ & $-6$ & $-4$ & $-5$ \\
& 0 & 1 & 2 & 1 & 2 \\
& 2 & 2 & 2 & 1 & 1 \\
\hline
41 & $-8$ & $-5$ & $-6$ & $-6$ & $-8$ \\
& 3 & 1 & 2 & 1 & 2 \\
& 2 & 2 & 2 & 1 & 1 \\
\hline
61 & $-10$ & $-9$ & $-12$ & $-8$ & $-10$ \\
& 3 & 2 & 2 & 3 & 2 \\
& 2 & 2 & 4 & 3 & 3 \\
\hline
71 & $-14$ & $-10$ & $-12$ & $-9$ & $-12$ \\
& 4 & 2 & 3 & 2 & 3 \\
& 2 & 3 & 5 & 2 & 2 \\
\end{tabular}
\qquad
\begin{tabular}{c|cccccc}
\multicolumn{7}{l}{Coefficients of} \\
$p$ & $x^5$ & $x^4$ & $x^3$ & $x^2$ & $x$ & $1 = 0$ \\
\hline
11 & 1 & 1 & $-4$ & $-3$ & $+3$ & $+1$ \\
31 & 1 & 1 & $-12$ & $-2$ & $+1$ & $+5$ \\
41 & 1 & 1 & $-16$ & $+5$ & $+21$ & $-9$ \\
61 & 1 & 1 & $-24$ & $-17$ & $+41$ & $-23$ \\
71 & 1 & 1 & $-28$ & $+37$ & $+25$ & $+1$ \\
\end{tabular}
\end{center}

\begin{center}
{\small 765]\hfill\hfill 317}
\end{center}

\begin{center}\textbf{765.}\end{center}

\begin{center}\textbf{ON THE FLEXURE AND EQUILIBRIUM OF A SKEW SURFACE.}\end{center}

\begin{center}\small [From the \textit{Proceedings of the London Mathematical Society}, vol.\ \textsc{xii}.\ (1881), pp.\ 103--108.\\
Read March 10, 1881.]\end{center}

\textsc{The} skew surface is taken to be such that the strip between two consecutive
generating lines is rigid, and that the flexure takes place by the rotation of the
strips about the generating lines successively. The theory of the flexure is well known,
but I am not aware that the theory of the equilibrium of such a surface, when acted
upon by any given forces, has been considered; it is, however, a question which
presents itself naturally in connexion with those relating to other continuous bodies
treated of in the \textit{M\'ecanique Analytique}, and forms a good example of the principles
made use of.

To begin with the mechanical theory: we may regard the forces as acting on
the generating lines regarded as material lines; and if for an element of mass $dm$,
coordinates $(x, y, z)$ of a particular generating line $G$, the forces parallel to the axes
are $X', Y', Z'$, then the corresponding term in the equation of equilibrium is
\[
S(X'\delta x + Y'\delta y + Z'\delta z)\,dm;
\]
and observing that there are (as will afterwards appear) five geometrical conditions, which
I represent by $U_1 = 0,\ U_2 = 0,\ \ldots,\ U_5 = 0$, the equation of equilibrium is
\[
S\{(X'\delta x + Y'\delta y + Z'\delta z)\,dm + T_1\delta U_1 + T_2\delta U_2 + T_3\delta U_3 + T_4\delta U_4 + T_5\delta U_5\} = 0,
\]
where $T_1, T_2, \ldots, T_5$ are the indeterminate multipliers, representing colligation-forces
which correspond to the five geometrical conditions respectively.

Taking $(\xi, \eta, \zeta)$ for the coordinates of a particular point $P$ on the generating
line; $p, q, r$ for the cos-inclinations of the line (whence $U_1 = p^2 + q^2 + r^2 - 1 = 0$ is one
of the geometrical relations), and $\rho$ for the distance of $dm$ from $P$, we have
\[
x, y, z = \xi + \rho p,\ \eta + \rho q,\ \zeta + \rho r,
\]
\[
\delta x, \delta y, \delta z = \delta\xi + \rho\delta p,\ \delta\eta + \rho\delta q,\ \delta\zeta + \rho\delta r.
\]

\begin{center}
{\small 318\hfill ON THE FLEXURE AND EQUILIBRIUM OF A SKEW SURFACE.\hfill [765}
\end{center}

The summation $S$ extends first to the different points of the generating line, and
then to the different generating lines; applying it first to the particular generating
line, we write
\[
\begin{array}{ccccccc}
SX'dm, & SY'dm, & SZ'dm, & SX'\rho\,dm, & SY'\rho\,dm, & SZ'\rho\,dm \\
= X, & Y, & Z, & L, & M, & N,
\end{array}
\]
where $X, Y, Z$ are the whole forces, and $L, M, N$ the whole moments about the
point $P$, for the generating line $G$; retaining the same summatory symbol $S$, as now
referring to the different generating lines, the equation becomes
\[
S\{X\delta\xi + Y\delta\eta + Z\delta\zeta + L\delta p + M\delta q + N\delta r + T_1\delta U_1 + \ldots + T_5\delta U_5\} = 0.
\]

We have now to consider the geometrical theory of the flexure. Taking on the
skew surface an arbitrary curve cutting each generating line $G$ in a point $P$, coordinates
$(\xi, \eta, \zeta)$, and taking $\sigma$ for the distance along the curve of the point $P$ from a fixed
point of the curve; also $p, q, r$, as before, for the cos-inclinations of the generating
line $G$, then when the surface is in a determinate state, $\xi, \eta, \zeta, p, q, r$ are given
functions of $\sigma$: but these functions vary with the flexure of the surface, with, however,
certain relations unaffected by the flexure; and the problem is to find first these
relations. As already mentioned, one of them is $p^2 + q^2 + r^2 - 1 = 0$.

Taking $P'$ as the consecutive point on the curve, so that the direction of the
element $PP'$ is that of the tangent $PT$ at $P$, it is convenient to write $l, m, n$ for
the cosine-inclinations of the tangent; we have, it is clear,
\[
l, m, n = \frac{d\xi}{d\sigma},\ \frac{d\eta}{d\sigma},\ \frac{d\zeta}{d\sigma};\qquad l^2 + m^2 + n^2 - 1 = 0.
\]
The conditions in order to the rigidity of the strip, are that the angles $GPP'$,
$G'P'P$ ($= 180^\circ - G'P'T$), and the inclination $G'P'$ to $GP$, shall have given values,
variable it may be from strip to strip---that is, these values must be given functions
of $\sigma$. Taking $\angle GPT = I$, the value of $G'P'T$ can differ only infinitesimally from that
of $GPT$, and we take it to be $G'P'T = I - \Omega d\sigma$; also the inclination $GP$ to $G'P'$
is an infinitesimal, $= \Theta d\sigma$: we have $I, \Omega, \Theta$ given functions of $\sigma$: It is to be
remarked that these conditions imply, inclination of $G'P'$ to tangent plane $GPT$ at $P$
has a given value $\Lambda d\sigma$; in fact, if through $P$ we draw a line $P\gamma$ parallel to $P'G'$,
then, if $P$ is regarded as the centre of a sphere which meets $PG, P\gamma, PT$ in the

\begin{center}
{\small 765]\hfill ON THE FLEXURE AND EQUILIBRIUM OF A SKEW SURFACE.\hfill 319}
\end{center}

points $g, g', t$ respectively, we have a spherical triangle $gg't$, the sides of which are
$I - \Omega d\sigma$, $I$, and $\Theta d\sigma$, and of which the perpendicular $g'm$ is $= \Lambda d\sigma$: we have thus

\textit{[Figure: small right-angled spherical triangle $gg't$ with base $\Omega d\sigma$, altitude $\Lambda d\sigma$, hypothenuse $\Theta d\sigma$, vertex $m$ on side $I$.]}

an infinitesimal right-angled triangle, the base and altitude of which are $\Omega d\sigma, \Lambda d\sigma$,
and the hypothenuse is $\Theta d\sigma$; whence $\Theta^2 = \Omega^2 + \Lambda^2$. In the case of the developable
surface $\Lambda = 0$ and $\Theta = \Omega$. It may be remarked that, when the curve on the skew
surface is the line of striction, we have $\Omega = 0$; in fact, taking $P$ to be on the line
of striction, the line
\[
\frac{X - \xi}{qr' - q'r} = \frac{Y - \eta}{rp' - r'p} = \frac{Z - \zeta}{pq' - p'q}
\]
through $(\xi, \eta, \zeta)$ at right angles to the two generating lines, meets the consecutive
generating line $X, Y, Z = \xi' + \rho p',\ \eta' + \rho q',\ \zeta' + \rho r'$; and the condition that this may
be so is easily found to be $\Omega = 0$.

Take, for a moment, $p', q', r'$ for the cos-inclinations of the consecutive generating
line $P'G'$; we have
\begin{align*}
lp + mq + nr &= \cos I, \\
lp' + mq' + nr' &= \cos(I - \Omega d\sigma), \\
pp' + qq' + rr' &= \cos\Theta d\sigma;
\end{align*}
and then writing $p', q', r' = p + dp,\ q + dq,\ r + dr$, and observing that the equation
$p'^2 + q'^2 + r'^2 = 1$ gives
\[
p\,dp + q\,dq + r\,dr = -\tfrac{1}{2}(dp^2 + dq^2 + dr^2),
\]
these equations and the before-mentioned two equations become
\begin{align*}
(U_1)\quad & p^2 + q^2 + r^2 - 1 = 0, \\
(U_2)\quad & l^2 + m^2 + n^2 - 1 = 0, \\
(U_3)\quad & lp + mq + nr - \cos I = 0, \\
(U_4)\quad & l\,dp + m\,dq + n\,dr - \Omega\sin I\,d\sigma = 0, \\
(U_5)\quad & dp^2 + dq^2 + dr^2 - \Theta^2 d\sigma^2 = 0,
\end{align*}
which equations, considering therein $l, m, n$ as standing for their values $\dfrac{d\xi}{d\sigma}, \dfrac{d\eta}{d\sigma}, \dfrac{d\zeta}{d\sigma}$,
are the geometrical relations which connect the six variables $\xi, \eta, \zeta, p, q, r$, considered
as functions of $\sigma$. And in these equations $I, \Omega, \Theta$ denote given functions of $\sigma$,
invariable by any flexure of the surface.

To complete the geometrical theory, it is to be observed that we can by flexure
bring the generating lines of the surface to be parallel to those of any given cone

\begin{center}
{\small 320\hfill ON THE FLEXURE AND EQUILIBRIUM OF A SKEW SURFACE.\hfill [765}
\end{center}

$C(p, q, r) = 0$, where $C(p, q, r)$ denotes a homogeneous function of $(p, q, r)$. Hence,
joining to the foregoing five equations this new equation
\[
C(p, q, r) = 0,
\]
these six equations determine $\xi, \eta, \zeta, p, q, r$ as functions of $\sigma$. To make the
solution completely determinate, we have only to assume for the point $P$, which
corresponds, say, to the value $\sigma = 0$, a position in space at pleasure, and to take the
corresponding generating line $PG$ parallel to a generating line, at pleasure, of the cone.

As an example, writing $\gamma$ to denote an arbitrary constant angle, if the invariable
conditions are
\[
I = \gamma,\ \Omega = \sin\gamma,\ \Theta = 0,
\]
then the five equations are
\begin{align*}
p^2 + q^2 + r^2 - 1 &= 0, \\
l^2 + m^2 + n^2 - 1 &= 0, \\
lp + mq + nr - \cos\gamma &= 0, \\
dp^2 + dq^2 + dr^2 - \sin^2\gamma\,d\sigma^2 &= 0, \\
l\,dp + m\,dq + n\,dr &= 0.
\end{align*}
We assume first
\[
C(p, q, r) = p^2 + q^2 - r^2\tan^2\gamma,\ = 0;
\]
and secondly
\[
C(p, q, r) = r,\ = 0.
\]

Then, in the former case, we find the solution
\[
p, q, r = -\sin\gamma\sin\sigma,\ \sin\gamma\cos\sigma,\ \cos\gamma
\]
\[
\xi, \eta, \zeta = \cos\sigma,\ \sin\sigma,\ 0;
\]
giving
\[
x, y, z = \cos\sigma - \rho\sin\gamma\sin\sigma,\ \sin\sigma + \rho\sin\gamma\cos\sigma,\ \rho\cos\gamma;
\]
and consequently
\[
x^2 + y^2 - \rho^2\tan^2\gamma = 0,
\]
the hyperboloid of revolution. And, in the latter case,
\[
p, q, r = \cos(\sigma\sin\gamma),\ \sin(\sigma\sin\gamma),\ 0,
\]
\[
\xi, \eta, \zeta = \cot\gamma\sin(\sigma\sin\gamma),\ -\cot\gamma\cos(\sigma\sin\gamma),\ \sigma\sin\gamma,
\]
that is,
\[
x, y = \cot\gamma\sin z + \rho\cos z,\ -\cot\gamma\cos z + \rho\sin z,
\]
whence
\[
x\sin z - y\cos z = \cot\gamma,
\]
a skew helicoid generated by horizontal tangents of the cylinder $x^2 + y^2 = \cot^2\gamma$. This
is a known deformation of the hyperboloid.

\begin{center}
{\small 765]\hfill ON THE FLEXURE AND EQUILIBRIUM OF A SKEW SURFACE.\hfill 321}
\end{center}

Returning now to the mechanical problem, we have to consider the terms
\begin{align*}
S\cdot\,& T_1\tfrac{1}{2}\delta(p^2 + q^2 + r^2 - 1) \\
+\,& T_2\tfrac{1}{2}\delta(l^2 + m^2 + n^2 - 1) \\
+\,& T_3\delta(lp + mq + nr - \cos I) \\
+\,& T_4\delta\left(l\frac{dp}{d\sigma} + m\frac{dq}{d\sigma} + n\frac{dr}{d\sigma} - \Omega\sin I\right) \\
+\,& T_5\tfrac{1}{2}\delta\left\{\left(\frac{dp}{d\sigma}\right)^2 + \left(\frac{dq}{d\sigma}\right)^2 + \left(\frac{dr}{d\sigma}\right)^2 - \Theta^2\right\}.
\end{align*}
The first term gives, under the sign $S$,
\[
T_1(p\delta p + q\delta q + r\delta r). \tag{*}
\]
The second term gives, in the first instance,
\[
\frac{T_2}{d\sigma}(l\,d\delta\xi + m\,d\delta\eta + n\,d\delta\zeta);
\]
or, since in general
\[
S\,l\,d\delta\xi = l^{(n)}\delta\xi^{(n)} - l^{(0)}\delta\xi^{(0)} + S^*(-dl\cdot\delta\xi),
\]
then, attending only to the terms under the sign $S^*$, these are
\[
= -\frac{d}{d\sigma}T_2 l\cdot\delta\xi - \frac{d}{d\sigma}T_2 m\cdot\delta\eta - \frac{d}{d\sigma}T_2 n\cdot\delta\zeta. \tag{*}
\]
The third term gives
\[
T_3(l\delta p + m\delta q + n\delta r) \tag{*}
\]
\[
+ T_3(p\delta l + q\delta m + r\delta n),
\]
where the second line,
\[
= \frac{T_3}{d\sigma}(p\,d\delta\xi + q\,d\delta\eta + r\,d\delta\zeta),
\]
attending only to the terms under the sign $S$, gives
\[
-\frac{d}{d\sigma}T_3 p\cdot\delta\xi - \frac{d}{d\sigma}T_3 q\cdot\delta\eta - \frac{d}{d\sigma}T_3 r\cdot\delta\zeta. \tag{*}
\]
The fourth term gives
\[
\frac{T_4}{d\sigma}\left(l\,d\delta p + m\,d\delta q + n\,d\delta r\right) + T_4\left(\frac{dp}{d\sigma}\delta l + \frac{dq}{d\sigma}\delta m + \frac{dr}{d\sigma}\delta n\right),
\]
where the first line, written under the form
\[
\frac{1}{d\sigma}\left\{d(T_4 l\,\delta p) - dT_4 l\cdot\delta p - T_4\,dl\cdot\delta p\right\}\ +\ \text{(similar for }m, n),
\]
and attending only to the terms under the sign $S$, gives
\[
-\frac{d}{d\sigma}T_4 l\cdot\delta p - \frac{d}{d\sigma}T_4 m\cdot\delta q - \frac{d}{d\sigma}T_4 n\cdot\delta r,
\]
\begin{center}{\small C. XI.\hfill 41}\end{center}

\begin{center}
{\small 322\hfill ON THE FLEXURE AND EQUILIBRIUM OF A SKEW SURFACE.\hfill [765}
\end{center}

and the second line, attending in like manner only to the terms under the sign $S$,
gives
\[
-\frac{d}{d\sigma}T_4 l\cdot\delta p - \frac{d}{d\sigma}T_4 m\cdot\delta q - \frac{d}{d\sigma}T_4 n\cdot\delta r. \tag{*}
\]

The fifth term, written under the form
\[
\frac{T_5}{d\sigma}\left(\frac{dp}{d\sigma}\,d\delta p + \frac{dq}{d\sigma}\,d\delta q + \frac{dr}{d\sigma}\,d\delta r\right),
\]
and attending only to the terms under the sign $S$, gives
\[
-\frac{d}{d\sigma}T_5\frac{dp}{d\sigma}\cdot\delta p - \frac{d}{d\sigma}T_5\frac{dq}{d\sigma}\cdot\delta q - \frac{d}{d\sigma}T_5\frac{dr}{d\sigma}\cdot\delta r; \tag{*}
\]
where in each case I have marked with an asterisk the lines which present them-selves
in the final result.

Hence, joining to the foregoing the force-terms
\[
X\delta\xi + Y\delta\eta + Z\delta\zeta + L\delta p + M\delta q + N\delta r, \tag{*}
\]
and equating to zero the coefficients of $\delta\xi, \delta\eta, \delta\zeta, \delta p, \delta q, \delta r$ respectively, we have
\[
\resizebox{\textwidth}{!}{$\displaystyle
\begin{array}{rcl}
0 &=& X - \dfrac{d}{d\sigma}T_2 l - \dfrac{d}{d\sigma}T_3 p, \\[2ex]
0 &=& Y - \dfrac{d}{d\sigma}T_2 m - \dfrac{d}{d\sigma}T_3 q, \\[2ex]
0 &=& Z - \dfrac{d}{d\sigma}T_2 n - \dfrac{d}{d\sigma}T_3 r, \\[2ex]
0 &=& L + T_1 p + T_3 l - \dfrac{d}{d\sigma}T_4 l - \dfrac{d}{d\sigma}T_5\dfrac{dp}{d\sigma}, \\[2ex]
0 &=& M + T_1 q + T_3 m - \dfrac{d}{d\sigma}T_4 m - \dfrac{d}{d\sigma}T_5\dfrac{dq}{d\sigma}, \\[2ex]
0 &=& N + T_1 r + T_3 n - \dfrac{d}{d\sigma}T_4 n - \dfrac{d}{d\sigma}T_5\dfrac{dr}{d\sigma},
\end{array}$}
\]
where it will be recollected that $l, m, n$ stand for $\dfrac{d\xi}{d\sigma}, \dfrac{d\eta}{d\sigma}, \dfrac{d\zeta}{d\sigma}$, the variables being
$\xi, \eta, \zeta, p, q, r$, and $\sigma$. The elimination of $T_1, T_2, \ldots, T_5$ from the six equations
should lead to a relation between $\xi, \eta, \zeta, p, q, r$, which, with the foregoing five
relations, would determine the six variables $\xi, \eta, \zeta, p, q, r$ in terms of $\sigma$.

In particular, the forces and moments $X, Y, Z, L, M, N$ may all of them
vanish; assuming that $T_1, T_2, \ldots, T_5$ do not all of them vanish, we still have the
sixth relation, which (with the foregoing five relations) determines $\xi, \eta, \zeta, p, q, r$ in
terms of $\sigma$; and it is to be remarked that the problem in question, of the figure
of equilibrium of the skew surface not acted upon by any forces, is analogous to
that of the geodesic line in space; only whilst here the solution is, curve a straight
line, the solution for the case of the skew surface depends upon equations of a
complex enough form; in the case of the developable surface, the required figure is
of course the plane.

\begin{center}
{\small 766]\hfill\hfill 323}
\end{center}

\begin{center}\textbf{766.}\end{center}

\begin{center}\textbf{ON THE GEODESIC CURVATURE OF A CURVE ON A SURFACE.}\end{center}

\begin{center}\small [From the \textit{Proceedings of the London Mathematical Society}, vol.\ \textsc{xii}.\ (1881), pp.\ 110--117.\\
Read April 14, 1881.]\end{center}

\textsc{There} is contained in Liouville's Note II.\ to his edition of Monge's \textit{Application
de l'Analyse \`a la G\'eom\'etrie} (Paris, 1850), see pp.\ 574 and 575, the following
formula,
\[
\frac{di}{ds} = -\frac{1}{2G\sqrt{E}}\frac{dG}{du}\cos i + \frac{1}{2E\sqrt{G}}\frac{dE}{dv}\sin i,
\]
\[
\frac{di}{ds} - \frac{\cos i}{\rho_u} + \frac{\sin i}{\rho_v},
\]
which gives the radius of geodesic curvature of a curve upon a surface when the
position of a point on the surface is defined by the parameters $u, v$, belonging to
a system of orthotomic curves; or, what is the same thing, such that
\[
ds^2 = E\,du^2 + G\,dv^2.
\]
Writing with Gauss $p, q$ instead of $u, v$, I propose to obtain the corresponding formula
in the general case where the parameters $p, q$ are such that
\[
ds^2 = E\,dp^2 + 2F\,dp\,dq + G\,dq^2.
\]

I call to mind that, if $PQ, PQ'$ are equal infinitesimal arcs on the given curve
and on its tangent geodesic, then the radius of geodesic curvature $\rho$ is, by definition,
a length $\rho$ such that $2\rho\cdot QQ' = PQ^2$. More generally, if the curves on the surface
are any two curves which touch each other, then $\rho$ as thus determined is the radius
of relative curvature of the two curves.

\begin{center}{\small 41---2}\end{center}

\begin{center}
{\small 324\hfill ON THE GEODESIC CURVATURE OF A CURVE ON A SURFACE.\hfill [766}
\end{center}

The notation is that of the Memoir, ``Disquisitiones generales circa superficies
curvas'' (1827), \textsc{Gauss}, \textit{Werke}, t.\ \textsc{iii}.; see also my paper ``On geodesic lines, in
particular those of a quadric surface,'' \textit{Proc.\ Lond.\ Math.\ Society}, t.\ \textsc{iv}.\ (1872),
pp.\ 191--211, [508]; and Salmon's \textit{Solid Geometry}, 3rd ed., 1874, pp.\ 251 et seq.

The coordinates $(x, y, z)$ of a point on the surface are taken to be functions of
two independent parameters $p, q$; and we then write
\begin{align*}
dx + \tfrac{1}{2}d^2x &= a\,dp + a'\,dq + \tfrac{1}{2}(\alpha\,dp^2 + 2\alpha'\,dp\,dq + \alpha''\,dq^2), \\
dy + \tfrac{1}{2}d^2y &= b\,dp + b'\,dq + \tfrac{1}{2}(\beta\,dp^2 + 2\beta'\,dp\,dq + \beta''\,dq^2), \\
dz + \tfrac{1}{2}d^2z &= c\,dp + c'\,dq + \tfrac{1}{2}(\gamma\,dp^2 + 2\gamma'\,dp\,dq + \gamma''\,dq^2):
\end{align*}
\[
E, F, G = a^2 + b^2 + c^2,\ aa' + bb' + cc',\ a'^2 + b'^2 + c'^2;\quad V^2 = EG - F^2;
\]
and therefore
\[
ds^2 = E\,dp^2 + 2F\,dp\,dq + G\,dq^2,
\]
where $E, F, G$ are regarded as given functions of $p$ and $q$.

To determine a curve on the surface, we establish a relation between the two
parameters $p, q$, or, what is the same thing, take $p, q$ to be functions of a single
parameter $\theta$; and we write as usual $p', p'', q'$, etc., to denote the differential
coefficients of $p, q$, etc., in regard to $\theta$; we write also $E_1, E_2$, etc., to denote the
differential coefficients $\dfrac{dE}{dp}, \dfrac{dE}{dq}$, etc. In the first instance, $\theta$ is taken to be an
arbitrary parameter, but we afterwards take it to be the length $s$ of the curve from
a fixed point thereof.

\begin{center}\textit{First formula for the radius of relative curvature.}\end{center}

Consider any two curves touching at the point $P$, coordinates $(x, y, z)$ which
are regarded as given functions of $(p, q)$; where $(p, q)$ are for the one curve given
functions, and for the other curve other given functions, of $\theta$.

The coordinates of a consecutive point for the one curve are then
\[
x + dx + \tfrac{1}{2}d^2x,\ y + dy + \tfrac{1}{2}d^2y,\ z + dz + \tfrac{1}{2}d^2z,
\]
where
\[
dp = p'\,d\theta + \tfrac{1}{2}p''\,d\theta^2,\quad dq = q'\,d\theta + \tfrac{1}{2}q''\,d\theta^2;
\]
hence these coordinates are
\[
x + (ap' + a'q')\,d\theta + \tfrac{1}{2}(\alpha p'^2 + 2\alpha' p'q' + \alpha'' q'^2)\,d\theta^2 + \tfrac{1}{2}(ap'' + a'q'')\,d\theta^2,
\]
and for the other curve they are in like manner
\[
x + (aP' + a'Q')\,d\theta + \tfrac{1}{2}(\alpha P'^2 + 2\alpha' P'Q' + \alpha'' Q'^2)\,d\theta^2 + \tfrac{1}{2}(aP'' + a'Q'')\,d\theta^2,
\]

\begin{center}
{\small 766]\hfill ON THE GEODESIC CURVATURE OF A CURVE ON A SURFACE.\hfill 325}
\end{center}

the only difference being in the terms which contain the second differential coefficients,
$p'', q''$ for the first curve, and $P'', Q''$ for the second curve. Hence the differences
of the coordinates are
\[
\tfrac{1}{2}\{a(p'' - P'') + a'(q'' - Q'')\}\,d\theta^2,\ \tfrac{1}{2}\{b(p'' - P'') + b'(q'' - Q'')\}\,d\theta^2,
\]
\[
\tfrac{1}{2}\{c(p'' - P'') + c'(q'' - Q'')\}\,d\theta^2,
\]
and consequently the distance $QQ'$ of the two consecutive points $Q, Q'$ is
\[
= \tfrac{1}{2}\sqrt{(E, F, G \mid p'' - P'', q'' - Q'')^2}\,d\theta^2.
\]
The squared arc $PQ^2$ is
\[
= (E, F, G \mid p', q')^2\,d\theta^2;
\]
and hence, if as before $2\rho\cdot QQ' = PQ^2$, that is, $\dfrac{1}{\rho} = 2QQ' \div PQ^2$, then
\[
\frac{1}{\rho} = \frac{\sqrt{(E, F, G \mid p'' - P'', q'' - Q'')^2}}{(E, F, G \mid p', q')^2},
\]
the required formula for $\rho$.

\begin{center}\textit{Second formula for the radius of relative curvature.}\end{center}

We now take the variable $\theta$ to be the length $s$ of the curve measured from a
fixed point thereof, so that $p', p''$, etc.\ denote $\dfrac{dp}{ds}, \dfrac{d^2p}{ds^2}$, etc. We have therefore
\[
1 = (E, F, G \mid p', q')^2,
\]
and the formula becomes
\[
\frac{1}{\rho} = \sqrt{(E, F, G \mid p'' - P'', q'' - Q'')^2}.
\]
But, differentiating the above equation as regards the curve, we find
\[
0 = 2(E, F, G \mid p', q')(p'', q'') + (E, F, G \mid p', q')^2,
\]
where $E, F, G$ are used to denote the complete differential coefficients $E_1 p' + E_2 q'$, etc.
And similarly, differentiating in regard to the tangent geodesic, we obtain
\[
0 = 2(E, F, G \mid p', q')(P'', Q'') + (E, F, G \mid p', q')^2;
\]
and hence, taking the difference of the two equations,
\[
0 = (E, F, G \mid p', q')(p'' - P'', q'' - Q'').
\]
Hence, in the equation for $\dfrac{1}{\rho}$, the function under the radical sign may be written
\[
(E, F, G \mid p', q')^2\cdot(E, F, G \mid p'' - P'', q'' - Q'')^2 - \{(E, F, G \mid p', q')(p'' - P'', q'' - Q'')\}^2,
\]

\begin{center}
{\small 326\hfill ON THE GEODESIC CURVATURE OF A CURVE ON A SURFACE.\hfill [766}
\end{center}

which is identically
\[
= (EG - F^2)\{p'(q'' - Q'') - q'(p'' - P'')\}^2.
\]
Hence, extracting the square root, and for $\sqrt{EG - F^2}$ writing $V$, we have
\[
\frac{1}{\rho} = V\{p'(q'' - Q'') - q'(p'' - P'')\},
\]
or say
\[
\frac{1}{\rho} = V(p'q'' - q'p'') - V(p'Q'' - q'P''),
\]
which is the new formula for the radius of relative curvature.

\begin{center}\textit{Formula for the radius of geodesic curvature.}\end{center}

In the paper ``On Geodesic Lines, etc.,'' p.\ 195, [vol.\ viii.\ of this Collection, p.\ 160],
writing $EG - F^2 = V^2$, and $P'', Q''$ in place of $p'', q''$, the differential equation of the
geodesic line is obtained in the form
\begin{align*}
& (Ep' + Fq')\{(2F_1 - E_2)p'^2 + 2G_1 p'q' + G_2 q'^2\} \\
& \quad - (Fp' + Gq')\{E_1 p'^2 + 2E_2 p'q' + (2F_2 - G_1)q'^2\} \\
& \quad\quad + 2V^2(p'Q'' - q'P'') = 0;
\end{align*}
or, denoting by $\Omega$ the first two lines of this equation, we have
\[
V(p'Q'' - q'P'') = -\tfrac{1}{2V}\Omega.
\]

The foregoing equation gives therefore, for the radius of geodesic curvature,
\[
\frac{1}{\rho} = V(p'q'' - p''q') + \tfrac{1}{2V}\Omega,
\]
which is an expression depending only upon $p', q'$, the first differential coefficients
(common to the curve and geodesic), and on $p'', q''$, the second differential coefficients
belonging to the curve.

Observe that $\Omega$ is a cubic function of $p', q'$: we have
\[
\Omega = (\mathfrak{A}, \mathfrak{B}, \mathfrak{C}, \mathfrak{D} \mid p', q')^3,
\]
the values of the coefficients being
\begin{align*}
\mathfrak{A} &= 2EF_1 - EE_2 - FE_1, \\
\mathfrak{B} &= 2EG_1 + 2FF_1 - 3FE_2 - GE_1, \\
\mathfrak{C} &= EG_2 + 2FG_1 - 2FF_2 - 2GE_2, \\
\mathfrak{D} &= FG_2 - 2GF_2 + GG_1.
\end{align*}

\end{document}
